\section{Graph-traversal packing}
\label{app:traversal}

This appendix specifies how documents are ordered and packed into a fixed-length
training sequence by traversing the inter-document link graph: the exact traversal
algorithms, the seed/budget bookkeeping, and --- most importantly --- the coupling
between \emph{traversal order} and the causal DAG gate of the cross-document mask.
The mask semantics themselves are given in \Cref{sec:method,app:kernels}; density
bucketing and the DDP scheduler that consume these packs are in \Cref{app:density}.

\subsection{Graph storage}
\label{app:traversal:graph}

The link graph is built once from the tokenized corpus. Each document is keyed by a
normalized identifier and assigned an integer id equal to its enumeration index in
JSONL insertion order; this order is a determinism dependency
(\Cref{app:traversal:det}). The graph is \textbf{directed}: for every node an
\emph{outgoing} neighbor list (documents it links to) and an \emph{incoming} list
(documents linking to it) are stored redundantly rather than derived, so no symmetry
is assumed. Adjacency is not deduplicated, so a repeated edge raises a random walk's
transition probability toward that neighbor (edge multiplicity acts as a weight). A
neighbor whose identifier is absent from the index is \textbf{silently dropped} at
construction time; consequently the \emph{realized} link density during training is
strictly below the raw graph density (\Cref{app:density}).

\subsection{Traversal strategies}
\label{app:traversal:strategies}

A strategy exposes two operations: reset (begin a fresh traversal from a seed) and
propose-next (return the next candidate document id, or signal frontier exhaustion).
Strategies produce only an \emph{order}; the token budget, per-pack deduplication, and
the stopping rule are enforced by the sampler (\Cref{app:traversal:budget}). The
default edge mode is \emph{outgoing}. Four strategies are wired into the training and
precompute paths; a fifth (a composite mixture) is present but unused in any reported
run.

\paragraph{BFS (deterministic).}
Breadth-first expansion from the seed using a FIFO frontier and a visited set.
\begin{enumerate}
  \item On reset, enqueue the seed and mark it visited.
  \item Propose-next pops the front node $u$; each not-yet-visited outgoing neighbor
    of $u$, taken \emph{in stored file order}, is marked visited on enqueue and
    appended to the frontier. Return $u$.
  \item On an empty frontier, restart from a uniformly random unvisited node (a bounded
    retry loop), which may jump to a disconnected component.
\end{enumerate}
Because neighbors are enqueued in a fixed file order, BFS consumes no randomness and is
fully deterministic given the seed.

\paragraph{DFS (stochastic).}
Depth-first expansion using a LIFO stack and a visited set. On reset the seed is pushed;
propose-next pops $u$, collects its unvisited neighbors, \emph{shuffles them with the
shared RNG} (for a canonical depth-first feel), pushes them marking visited, and returns
$u$. Empty-stack restart is as in BFS. Unlike BFS, DFS \emph{consumes the shared RNG} on
every expansion; this asymmetry shifts the shared random stream and matters for
reproducibility (\Cref{app:traversal:det}).

\paragraph{RandomWalk (first-order Markov).}
A single current node with no visited set, so revisits are allowed. Each step teleports
with probability $p_\text{restart}$ to a \textbf{uniformly random node over the whole
graph}; otherwise it follows an edge. Under edge mode \emph{both}, a direction is first
drawn with probability proportional to weights $(w_\text{in}, w_\text{out})$ and then a
neighbor uniformly within that side; edge modes \emph{in}/\emph{out} use only the
corresponding side. A dead end (no eligible neighbor) forces a teleport. The chain is
first-order Markov with stationary distribution proportional to out-degree, carrying no
PageRank damping or normalization. \textbf{Because the restart teleports to a uniform
node rather than back to the seed, this is not Random Walk with Restart or personalized
PageRank} \citep{page1999pagerank,tong2006rwr}; those restart to a fixed seed or
personalization vector and sample a different distribution.

\paragraph{Random (structure-free control).}
Ignores the graph: every proposal draws a document id uniformly over all ids. A
``Random'' pack is $N$ uniformly sampled documents and serves as the no-structure
baseline.

\subsection{Seed selection, token budget, and pack assembly}
\label{app:traversal:budget}

\paragraph{Uniform seeding, independent of strategy.}
A pack is grown from one or more \emph{subwalks}, and every subwalk seed is drawn by
\textbf{uniform rejection sampling over all documents regardless of the traversal
strategy}; the strategy governs only how the neighborhood grows from that seed. A
``BFS packing'' run is therefore a set of uniform-seeded local BFS neighborhoods, not a
single global sweep, and likewise for DFS and RandomWalk --- a deliberate control (every
document is equally likely to anchor a pack) that nonetheless shapes what each condition
measures.

\paragraph{Token budget and per-document cap.}
The budget of a pack is $\texttt{batch\_size}\times\texttt{seq\_len}$ (the maximum
sequence length), charged against the \emph{full} per-document layout: prefix
identifier-card $+$ body $+$ EOS suffix (\Cref{app:density}). A per-document body cap
bounds how many body tokens any one document contributes; an over-cap document is either
body-truncated or skipped, per the overflow policy.

\paragraph{Subwalks per pack.}
A \emph{fresh} strategy object is instantiated per subwalk and grown until the budget is
(nearly) reached or no valid seed remains, then another subwalk begins, so one pack can
hold several disconnected neighborhoods. Reset is thus invoked per \emph{subwalk}, not
per pack.

\paragraph{Exact-length truncation.}
Assembly may overshoot the budget by at most one document; the pack is then trimmed to
\emph{exactly} the target length. Exactness is not cosmetic: the Triton attention
kernels take the sequence length as a compile-time constant, so a variable length would
trigger a full re-autotune (order $140$\,s as-measured) and desynchronize the DDP ranks.
Trimming removes \emph{body} tokens only, never prefix/EOS decoration, walking documents
from the head by default and carrying any residual forward. Under the default
targets-first ordering below, head trimming sheds tokens from the earliest documents,
which are exactly the link targets.

\subsection{Order mode and the traversal--DAG-gate coupling}
\label{app:traversal:order}

The layout order of the collected documents is set by one of two modes:
\emph{as-traversed} (proposal order) or \emph{prefer-targets-first}, the
\textbf{default}. This choice is the single most important non-obvious design point in
the pipeline, because it is forced by the cross-document mask.

\paragraph{Why targets must precede linkers.}
The mask is causal at its base (a query at position $q$ attends only to keys at
$k\le q$), and a detected link $A\!\to\!B$ opens a grant letting linker $A$ attend into
target $B$ only if $B$ begins before the link position; otherwise a DAG gate silently
drops the grant, never relaxing causality (\Cref{app:kernels}). With outgoing traversal
the natural proposal order places linker $A$ before target $B$, so under as-traversed the
target lands \emph{after} the linker, the gate fires, and the edge is never realized.
Hence:
\begin{center}
\emph{with outgoing traversal, a targets-before-linkers reordering is required for the
causal cross-document mask to realize any link edge at all.}
\end{center}

\paragraph{Prefer-targets-first: per-component topological sort.}
The default runs Kahn's algorithm \emph{within each connected component}, emitting every
link's target before its linker, and lays each component down as a \textbf{contiguous
block}. Contiguity is a hard requirement of the concatenation kernel, which merges a
component only if its positions form a single interval (\Cref{app:kernels}), and it is
what makes the analytic density proxy of \Cref{app:density} exact, since every target
then ends before its link position. Components are \emph{structural}: undirected
union-find components over the graph edges, deliberately decoupled from traversal order
so that a frontier-exhaustion restart pulling in an unrelated document does not merge two
unrelated repositories into one attention block. Within a component containing a
\emph{cycle}, Kahn's algorithm cannot fully order the nodes and falls back to insertion
order for the offending ones; links then ordered against the gate are dropped, so cyclic
link structure is realized only partially.

\subsection{Determinism and a documented inconsistency}
\label{app:traversal:det}

A single shared pseudo-random generator, seeded per rank as $\texttt{base\_seed}+
\texttt{rank}$, drives \emph{both} seed selection and every proposal. Because the stream
is shared, any strategy that consumes randomness (DFS shuffling, RandomWalk teleports)
shifts all subsequent draws, making a run \emph{stream-order fragile}: two configurations
differing only in a strategy's RNG consumption diverge everywhere downstream. Determinism
holds only with synchronous single-stream data loading (zero worker processes) and
depends on the JSONL insertion order defining the document ids.

The live training path and the offline epoch-precompute path use different RandomWalk
restart probabilities --- $p_\text{restart}=0.05$ live versus $0.0$ in precompute --- so
they sample subtly different walk distributions, and any comparison mixing live and
precomputed RandomWalk packs inherits this discrepancy.

\subsection{Relation to node-sampling and random-walk literature}
\label{app:traversal:related}

Traversal-ordered packing is closest in spirit to node-sampling in graph representation
learning, but the objective differs: we order documents so a causal attention mask can
realize link edges, not to learn node embeddings. The RandomWalk strategy is a plain
first-order walk, like the context-generating walks of DeepWalk
\citep{perozzi2014deepwalk}, but without the second-order $p,q$ in/out bias of node2vec
\citep{grover2016node2vec}; the edge-mode weights give only a first-order directional
preference. BFS and DFS mirror the breadth- and depth-oriented neighborhood samplers
that node2vec interpolates between, and the uniform-seeded, budget-capped local
neighborhoods resemble the fixed-size neighbor sampling of GraphSAGE
\citep{hamilton2017graphsage}. Against the PageRank family, the uniform-teleport restart
means the walk implements neither personalized PageRank nor Random Walk with Restart
\citep{page1999pagerank,tong2006rwr}. The novel element is not the walk but the
coordinated targets-first topological reordering that lets a \emph{causal} cross-document
mask realize link edges (\Cref{app:traversal:order}), together with structural union-find
components robust to frontier restarts and the exact-length packing forced by the
kernel's compile-time sequence length.
